% --- Archon physics formalization source begin ---
% archon:physics
% archon:covers PhyXMiniProblems/problem_phyx_mini_0166.lean
% archon:source-report reports/phyx_mini/problem_phyx_mini_0166.source.json

\chapter{Physics problem phyx\_mini\_0166}
\label{ch:PhyXMiniProblems_problem_phyx_mini_0166}

\paragraph{Problem source.}
\#\# Physical scenario

A vibrating string \textbackslash{}( 50.0\textbackslash{},\textbackslash{}text\{cm\} \textbackslash{}) long is under a tension of \textbackslash{}( 1.00\textbackslash{},\textbackslash{}text\{N\} \textbackslash{}). The results from five successive stroboscopic pictures are shown in \textbackslash{}textbf\{figure\}. The strobe rate is set at \textbackslash{}( 5000 \textbackslash{}) flashes per minute, and observations reveal that the maximum displacement occurred at flashes 1 and 5 with no other maxima in between.

\#\# Question

What is the mass of this string?

\#\# Answer choices

- A: A: 18.5g

- B: B: 16.5g

- C: C: 19.5g

- D: D: 48.5g

\#\# Figure caption (auxiliary; use the image as primary evidence)

The image depicts a diagram illustrating wave patterns, likely related to standing waves on a string between two fixed points.

- **Wave Patterns**: The diagram shows several overlapping wave curves, each labeled with numbers (1 to 5). 
- **Colors**: Each wave is a different color: black, green, blue, and magenta.
- **Nodes and Antinodes**: The waves intersect at the midpoint, indicating nodes where there is no displacement.
- **Labels**:
  - A point on the black wave is labeled "P".
  - Each wave is numbered at its peak positions, from 1 to 5.
- **Measurements**: The vertical lines at one end indicate a total height of 3 cm, divided into two sections of 1.5 cm each.
- **Boundaries**: Vertical lines on the left and right suggest the fixed endpoints of the string or medium.

This setup visually represents the concept of standing waves, showing how waves of different modes fit within the fixed length, with nodes at the endpoints and antinodes at the peaks.

\#\# Dataset metadata

Resonance and Harmonics; Physical Model Grounding Reasoning, Multi-Formula Reasoning

\paragraph{Recorded answer/context.}
A: A: 18.5g

\paragraph{Figure/image path.}
/root/proposal\_for\_physic/science-mango/phyx\_mini\_run/phyx\_data/test\_image/166.png

\paragraph{Formalization target.}
create a compiling Lean file with sorry bodies at `PhyXMiniProblems/problem\_phyx\_mini\_0166.lean`. The Lean declarations must preserve the physical quantities, dimensions or dimensional roles, figure labels, governing-law hypotheses, and final relation expressed by this problem.
Use Mathlib/Physlib names found through LeanExplore where available. If a physics API is missing, introduce faithful local abstractions rather than scalar placeholder aliases.

\begin{theorem}[Physics formalization target]
\label{thm:physics:phyx_mini_0166:target}
\lean{PhyXMiniProblems.ProblemPhyXMini0166.problem_phyx_mini_0166}
\uses{def:physics:phyx-mini-0166:phyxminiproblems-problemphyxmini0166-massingrams, def:physics:phyx-mini-0166:phyxminiproblems-problemphyxmini0166-vibratingstringsetup, def:physics:phyx-mini-0166:phyxminiproblems-problemphyxmini0166-matchesproblemreadouts, def:physics:phyx-mini-0166:phyxminiproblems-problemphyxmini0166-matchesprimaryfigure, def:physics:phyx-mini-0166:phyxminiproblems-problemphyxmini0166-haspositivestringquantities, def:physics:phyx-mini-0166:phyxminiproblems-problemphyxmini0166-obeysstandingstringlaws, def:physics:phyx-mini-0166:phyxminiproblems-problemphyxmini0166-isclosestdisplayedmass}
The assigned autoformalize agent should translate this physics problem into Lean declarations in the covered file, with theorem and lemma proof bodies written as `by sorry`.
\end{theorem}
\begin{proof}
This is an autoformalization task, not a proof task. Produce statements that can later be proved by the physics prover without weakening the physical model.
\end{proof}

% --- Archon declaration topology begin ---
\section{Declaration topology}
\begin{definition}[Physical Length]
  \label{def:physics:phyx-mini-0166:phyxminiproblems-problemphyxmini0166-physicallength}
  \lean{PhyXMiniProblems.ProblemPhyXMini0166.PhysicalLength}
  A nonnegative physical length, independent of the selected length unit
\end{definition}
\begin{proof}
  This is the declared model definition of Physical Length.
\end{proof}
\begin{definition}[Transverse Displacement]
  \label{def:physics:phyx-mini-0166:phyxminiproblems-problemphyxmini0166-transversedisplacement}
  \lean{PhyXMiniProblems.ProblemPhyXMini0166.TransverseDisplacement}
  A signed transverse displacement, independent of the selected length unit
\end{definition}
\begin{proof}
  This is the declared model definition of Transverse Displacement.
\end{proof}
\begin{definition}[Physical Time]
  \label{def:physics:phyx-mini-0166:phyxminiproblems-problemphyxmini0166-physicaltime}
  \lean{PhyXMiniProblems.ProblemPhyXMini0166.PhysicalTime}
  A nonnegative physical time interval
\end{definition}
\begin{proof}
  This is the declared model definition of Physical Time.
\end{proof}
\begin{definition}[Physical Frequency]
  \label{def:physics:phyx-mini-0166:phyxminiproblems-problemphyxmini0166-physicalfrequency}
  \lean{PhyXMiniProblems.ProblemPhyXMini0166.PhysicalFrequency}
  A nonnegative physical frequency
\end{definition}
\begin{proof}
  This is the declared model definition of Physical Frequency.
\end{proof}
\begin{definition}[Physical Mass]
  \label{def:physics:phyx-mini-0166:phyxminiproblems-problemphyxmini0166-physicalmass}
  \lean{PhyXMiniProblems.ProblemPhyXMini0166.PhysicalMass}
  A nonnegative physical mass
\end{definition}
\begin{proof}
  This is the declared model definition of Physical Mass.
\end{proof}
\begin{definition}[Physical Force]
  \label{def:physics:phyx-mini-0166:phyxminiproblems-problemphyxmini0166-physicalforce}
  \lean{PhyXMiniProblems.ProblemPhyXMini0166.PhysicalForce}
  A nonnegative force; its SI scalar readout is measured in newtons
\end{definition}
\begin{proof}
  This is the declared model definition of Physical Force.
\end{proof}
\begin{definition}[Linear Mass Density]
  \label{def:physics:phyx-mini-0166:phyxminiproblems-problemphyxmini0166-linearmassdensity}
  \lean{PhyXMiniProblems.ProblemPhyXMini0166.LinearMassDensity}
  Mass per unit length of the uniform string
\end{definition}
\begin{proof}
  This is the declared model definition of Linear Mass Density.
\end{proof}
\begin{definition}[centimeter Unit Choices]
  \label{def:physics:phyx-mini-0166:phyxminiproblems-problemphyxmini0166-centimeterunitchoices}
  \lean{PhyXMiniProblems.ProblemPhyXMini0166.centimeterUnitChoices}
  SI base units with centimeters selected as the length unit
\end{definition}
\begin{proof}
  This is the declared model definition of centimeter Unit Choices.
\end{proof}
\begin{definition}[gram Unit Choices]
  \label{def:physics:phyx-mini-0166:phyxminiproblems-problemphyxmini0166-gramunitchoices}
  \lean{PhyXMiniProblems.ProblemPhyXMini0166.gramUnitChoices}
  SI base units with grams selected as the mass unit
\end{definition}
\begin{proof}
  This is the declared model definition of gram Unit Choices.
\end{proof}
\begin{definition}[minute Unit Choices]
  \label{def:physics:phyx-mini-0166:phyxminiproblems-problemphyxmini0166-minuteunitchoices}
  \lean{PhyXMiniProblems.ProblemPhyXMini0166.minuteUnitChoices}
  SI base units with minutes selected as the time unit
\end{definition}
\begin{proof}
  This is the declared model definition of minute Unit Choices.
\end{proof}
\begin{definition}[length In Centimeters]
  \label{def:physics:phyx-mini-0166:phyxminiproblems-problemphyxmini0166-lengthincentimeters}
  \lean{PhyXMiniProblems.ProblemPhyXMini0166.lengthInCentimeters}
  \uses{def:physics:phyx-mini-0166:phyxminiproblems-problemphyxmini0166-physicallength, def:physics:phyx-mini-0166:phyxminiproblems-problemphyxmini0166-centimeterunitchoices}
  Scalar readout of a physical length in centimeters
\end{definition}
\begin{proof}
  This is the declared model definition of length In Centimeters.
\end{proof}
\begin{definition}[length In Meters]
  \label{def:physics:phyx-mini-0166:phyxminiproblems-problemphyxmini0166-lengthinmeters}
  \lean{PhyXMiniProblems.ProblemPhyXMini0166.lengthInMeters}
  \uses{def:physics:phyx-mini-0166:phyxminiproblems-problemphyxmini0166-physicallength}
  Scalar readout of a physical length in meters
\end{definition}
\begin{proof}
  This is the declared model definition of length In Meters.
\end{proof}
\begin{definition}[displacement In Centimeters]
  \label{def:physics:phyx-mini-0166:phyxminiproblems-problemphyxmini0166-displacementincentimeters}
  \lean{PhyXMiniProblems.ProblemPhyXMini0166.displacementInCentimeters}
  \uses{def:physics:phyx-mini-0166:phyxminiproblems-problemphyxmini0166-transversedisplacement, def:physics:phyx-mini-0166:phyxminiproblems-problemphyxmini0166-centimeterunitchoices}
  Scalar readout of a signed transverse displacement in centimeters
\end{definition}
\begin{proof}
  This is the declared model definition of displacement In Centimeters.
\end{proof}
\begin{definition}[time In Seconds]
  \label{def:physics:phyx-mini-0166:phyxminiproblems-problemphyxmini0166-timeinseconds}
  \lean{PhyXMiniProblems.ProblemPhyXMini0166.timeInSeconds}
  \uses{def:physics:phyx-mini-0166:phyxminiproblems-problemphyxmini0166-physicaltime}
  Scalar readout of a physical time in seconds
\end{definition}
\begin{proof}
  This is the declared model definition of time In Seconds.
\end{proof}
\begin{definition}[frequency In Hertz]
  \label{def:physics:phyx-mini-0166:phyxminiproblems-problemphyxmini0166-frequencyinhertz}
  \lean{PhyXMiniProblems.ProblemPhyXMini0166.frequencyInHertz}
  \uses{def:physics:phyx-mini-0166:phyxminiproblems-problemphyxmini0166-physicalfrequency}
  Scalar readout of a physical frequency in hertz
\end{definition}
\begin{proof}
  This is the declared model definition of frequency In Hertz.
\end{proof}
\begin{definition}[frequency Per Minute]
  \label{def:physics:phyx-mini-0166:phyxminiproblems-problemphyxmini0166-frequencyperminute}
  \lean{PhyXMiniProblems.ProblemPhyXMini0166.frequencyPerMinute}
  \uses{def:physics:phyx-mini-0166:phyxminiproblems-problemphyxmini0166-physicalfrequency, def:physics:phyx-mini-0166:phyxminiproblems-problemphyxmini0166-minuteunitchoices}
  Scalar readout of a physical frequency per minute
\end{definition}
\begin{proof}
  This is the declared model definition of frequency Per Minute.
\end{proof}
\begin{definition}[force In Newtons]
  \label{def:physics:phyx-mini-0166:phyxminiproblems-problemphyxmini0166-forceinnewtons}
  \lean{PhyXMiniProblems.ProblemPhyXMini0166.forceInNewtons}
  \uses{def:physics:phyx-mini-0166:phyxminiproblems-problemphyxmini0166-physicalforce}
  Scalar SI readout of a force in newtons
\end{definition}
\begin{proof}
  This is the declared model definition of force In Newtons.
\end{proof}
\begin{definition}[mass In Grams]
  \label{def:physics:phyx-mini-0166:phyxminiproblems-problemphyxmini0166-massingrams}
  \lean{PhyXMiniProblems.ProblemPhyXMini0166.massInGrams}
  \uses{def:physics:phyx-mini-0166:phyxminiproblems-problemphyxmini0166-physicalmass, def:physics:phyx-mini-0166:phyxminiproblems-problemphyxmini0166-gramunitchoices}
  Scalar readout of a mass in grams
\end{definition}
\begin{proof}
  This is the declared model definition of mass In Grams.
\end{proof}
\begin{definition}[Flash Label]
  \label{def:physics:phyx-mini-0166:phyxminiproblems-problemphyxmini0166-flashlabel}
  \lean{PhyXMiniProblems.ProblemPhyXMini0166.FlashLabel}
  The labels printed beside the five successive stroboscopic profiles
\end{definition}
\begin{proof}
  This is the declared model definition of Flash Label.
\end{proof}
\begin{definition}[ordinal]
  \label{def:physics:phyx-mini-0166:phyxminiproblems-problemphyxmini0166-flashlabel-ordinal}
  \lean{PhyXMiniProblems.ProblemPhyXMini0166.FlashLabel.ordinal}
  \uses{def:physics:phyx-mini-0166:phyxminiproblems-problemphyxmini0166-flashlabel}
  The one-based ordinal printed next to a stroboscopic profile
\end{definition}
\begin{proof}
  This is the declared model definition of ordinal.
\end{proof}
\begin{definition}[Figure Point]
  \label{def:physics:phyx-mini-0166:phyxminiproblems-problemphyxmini0166-figurepoint}
  \lean{PhyXMiniProblems.ProblemPhyXMini0166.FigurePoint}
  Geometrically distinguished points in the primary standing-wave figure
\end{definition}
\begin{proof}
  This is the declared model definition of Figure Point.
\end{proof}
\begin{definition}[Vibrating String Setup]
  \label{def:physics:phyx-mini-0166:phyxminiproblems-problemphyxmini0166-vibratingstringsetup}
  \lean{PhyXMiniProblems.ProblemPhyXMini0166.VibratingStringSetup}
  \uses{def:physics:phyx-mini-0166:phyxminiproblems-problemphyxmini0166-physicallength, def:physics:phyx-mini-0166:phyxminiproblems-problemphyxmini0166-transversedisplacement, def:physics:phyx-mini-0166:phyxminiproblems-problemphyxmini0166-physicaltime, def:physics:phyx-mini-0166:phyxminiproblems-problemphyxmini0166-physicalfrequency, def:physics:phyx-mini-0166:phyxminiproblems-problemphyxmini0166-physicalmass, def:physics:phyx-mini-0166:phyxminiproblems-problemphyxmini0166-physicalforce, def:physics:phyx-mini-0166:phyxminiproblems-problemphyxmini0166-linearmassdensity, def:physics:phyx-mini-0166:phyxminiproblems-problemphyxmini0166-flashlabel, def:physics:phyx-mini-0166:phyxminiproblems-problemphyxmini0166-figurepoint}
  All physical quantities required by the calculation and all labeled profile readouts from the primary figure. modeNumber is the number of half wavelengths along the fixed string
\end{definition}
\begin{proof}
  This is the declared model definition of Vibrating String Setup.
\end{proof}
\begin{definition}[Matches Problem Readouts]
  \label{def:physics:phyx-mini-0166:phyxminiproblems-problemphyxmini0166-matchesproblemreadouts}
  \lean{PhyXMiniProblems.ProblemPhyXMini0166.MatchesProblemReadouts}
  \uses{def:physics:phyx-mini-0166:phyxminiproblems-problemphyxmini0166-lengthincentimeters, def:physics:phyx-mini-0166:phyxminiproblems-problemphyxmini0166-frequencyperminute, def:physics:phyx-mini-0166:phyxminiproblems-problemphyxmini0166-forceinnewtons, def:physics:phyx-mini-0166:phyxminiproblems-problemphyxmini0166-vibratingstringsetup}
  Numerical measurements stated in the text and printed on the figure. In particular, no mass, density, wave speed, wavelength, period, or oscillation frequency is specified here
\end{definition}
\begin{proof}
  This is the declared model definition of Matches Problem Readouts.
\end{proof}
\begin{definition}[Matches Primary Figure]
  \label{def:physics:phyx-mini-0166:phyxminiproblems-problemphyxmini0166-matchesprimaryfigure}
  \lean{PhyXMiniProblems.ProblemPhyXMini0166.MatchesPrimaryFigure}
  \uses{def:physics:phyx-mini-0166:phyxminiproblems-problemphyxmini0166-lengthincentimeters, def:physics:phyx-mini-0166:phyxminiproblems-problemphyxmini0166-displacementincentimeters, def:physics:phyx-mini-0166:phyxminiproblems-problemphyxmini0166-vibratingstringsetup}
  Readouts and qualitative geometry taken from the primary figure. The two lobes have opposite transverse displacements at every flash. Flashes 1 and 5 are opposite extrema at P, and none of the intervening flashes is extremal
\end{definition}
\begin{proof}
  This is the declared model definition of Matches Primary Figure.
\end{proof}
\begin{definition}[Has Positive String Quantities]
  \label{def:physics:phyx-mini-0166:phyxminiproblems-problemphyxmini0166-haspositivestringquantities}
  \lean{PhyXMiniProblems.ProblemPhyXMini0166.HasPositiveStringQuantities}
  \uses{def:physics:phyx-mini-0166:phyxminiproblems-problemphyxmini0166-lengthincentimeters, def:physics:phyx-mini-0166:phyxminiproblems-problemphyxmini0166-lengthinmeters, def:physics:phyx-mini-0166:phyxminiproblems-problemphyxmini0166-timeinseconds, def:physics:phyx-mini-0166:phyxminiproblems-problemphyxmini0166-frequencyinhertz, def:physics:phyx-mini-0166:phyxminiproblems-problemphyxmini0166-forceinnewtons, def:physics:phyx-mini-0166:phyxminiproblems-problemphyxmini0166-massingrams, def:physics:phyx-mini-0166:phyxminiproblems-problemphyxmini0166-vibratingstringsetup}
  Strict positivity of the physical magnitudes used in the wave laws
\end{definition}
\begin{proof}
  This is the declared model definition of Has Positive String Quantities.
\end{proof}
\begin{definition}[Obeys Standing String Laws]
  \label{def:physics:phyx-mini-0166:phyxminiproblems-problemphyxmini0166-obeysstandingstringlaws}
  \lean{PhyXMiniProblems.ProblemPhyXMini0166.ObeysStandingStringLaws}
  \uses{def:physics:phyx-mini-0166:phyxminiproblems-problemphyxmini0166-flashlabel, def:physics:phyx-mini-0166:phyxminiproblems-problemphyxmini0166-vibratingstringsetup}
  The physical laws used in the calculation, stated in every coherent choice of units: flash rate and flash interval are reciprocal; consecutive opposite extrema are separated by half a period (four flash intervals in this observation); oscillation frequency and period are reciprocal; a fixed-string mode contains modeNumber half wavelengths; wave speed is wavelength times frequency; the stretched-string law is v² μ = tension; the string is uniform, so mass = μ length.
\end{definition}
\begin{proof}
  This is the declared model definition of Obeys Standing String Laws.
\end{proof}
\begin{definition}[Answer Choice]
  \label{def:physics:phyx-mini-0166:phyxminiproblems-problemphyxmini0166-answerchoice}
  \lean{PhyXMiniProblems.ProblemPhyXMini0166.AnswerChoice}
  Labels of the four displayed multiple-choice answers
\end{definition}
\begin{proof}
  This is the declared model definition of Answer Choice.
\end{proof}
\begin{definition}[mass In Grams]
  \label{def:physics:phyx-mini-0166:phyxminiproblems-problemphyxmini0166-answerchoice-massingrams}
  \lean{PhyXMiniProblems.ProblemPhyXMini0166.AnswerChoice.massInGrams}
  \uses{def:physics:phyx-mini-0166:phyxminiproblems-problemphyxmini0166-massingrams, def:physics:phyx-mini-0166:phyxminiproblems-problemphyxmini0166-answerchoice}
  The mass readout in grams printed beside an answer choice
\end{definition}
\begin{proof}
  This is the declared model definition of mass In Grams.
\end{proof}
\begin{definition}[Is Closest Displayed Mass]
  \label{def:physics:phyx-mini-0166:phyxminiproblems-problemphyxmini0166-isclosestdisplayedmass}
  \lean{PhyXMiniProblems.ProblemPhyXMini0166.IsClosestDisplayedMass}
  \uses{def:physics:phyx-mini-0166:phyxminiproblems-problemphyxmini0166-massingrams, def:physics:phyx-mini-0166:phyxminiproblems-problemphyxmini0166-vibratingstringsetup, def:physics:phyx-mini-0166:phyxminiproblems-problemphyxmini0166-answerchoice}
  A displayed choice is at least as close as every other displayed mass
\end{definition}
\begin{proof}
  This is the declared model definition of Is Closest Displayed Mass.
\end{proof}
\begin{lemma}[oscillation Period In Seconds eq twelve div one Hundred Twenty Five]
  \label{lem:physics:phyx-mini-0166:phyxminiproblems-problemphyxmini0166-oscillationperiodinseconds-eq-twelve-div-onehundredtwentyfive}
  \lean{PhyXMiniProblems.ProblemPhyXMini0166.oscillationPeriodInSeconds_eq_twelve_div_oneHundredTwentyFive}
  \uses{def:physics:phyx-mini-0166:phyxminiproblems-problemphyxmini0166-timeinseconds, def:physics:phyx-mini-0166:phyxminiproblems-problemphyxmini0166-vibratingstringsetup, def:physics:phyx-mini-0166:phyxminiproblems-problemphyxmini0166-matchesproblemreadouts, def:physics:phyx-mini-0166:phyxminiproblems-problemphyxmini0166-haspositivestringquantities, def:physics:phyx-mini-0166:phyxminiproblems-problemphyxmini0166-obeysstandingstringlaws}
  The four flash intervals from flash 1 to the opposite extremum at flash 5 constitute half a cycle. At 5000 flashes per minute, the full period is 12/125 s
\end{lemma}
\begin{proof}
  The conclusion follows from the governing relations and figure readouts named in the statement of oscillation Period In Seconds eq twelve div one Hundred Twenty Five.
\end{proof}
\begin{lemma}[oscillation Frequency In Hertz eq one Hundred Twenty Five div twelve]
  \label{lem:physics:phyx-mini-0166:phyxminiproblems-problemphyxmini0166-oscillationfrequencyinhertz-eq-onehundredtwentyfive-div-twelve}
  \lean{PhyXMiniProblems.ProblemPhyXMini0166.oscillationFrequencyInHertz_eq_oneHundredTwentyFive_div_twelve}
  \uses{def:physics:phyx-mini-0166:phyxminiproblems-problemphyxmini0166-frequencyinhertz, def:physics:phyx-mini-0166:phyxminiproblems-problemphyxmini0166-vibratingstringsetup, def:physics:phyx-mini-0166:phyxminiproblems-problemphyxmini0166-matchesproblemreadouts, def:physics:phyx-mini-0166:phyxminiproblems-problemphyxmini0166-haspositivestringquantities, def:physics:phyx-mini-0166:phyxminiproblems-problemphyxmini0166-obeysstandingstringlaws}
  The corresponding oscillation frequency is 125/12 Hz
\end{lemma}
\begin{proof}
  The conclusion follows from the governing relations and figure readouts named in the statement of oscillation Frequency In Hertz eq one Hundred Twenty Five div twelve.
\end{proof}
\begin{lemma}[wavelength In Meters eq one Half]
  \label{lem:physics:phyx-mini-0166:phyxminiproblems-problemphyxmini0166-wavelengthinmeters-eq-onehalf}
  \lean{PhyXMiniProblems.ProblemPhyXMini0166.wavelengthInMeters_eq_oneHalf}
  \uses{def:physics:phyx-mini-0166:phyxminiproblems-problemphyxmini0166-lengthinmeters, def:physics:phyx-mini-0166:phyxminiproblems-problemphyxmini0166-vibratingstringsetup, def:physics:phyx-mini-0166:phyxminiproblems-problemphyxmini0166-matchesproblemreadouts, def:physics:phyx-mini-0166:phyxminiproblems-problemphyxmini0166-matchesprimaryfigure, def:physics:phyx-mini-0166:phyxminiproblems-problemphyxmini0166-haspositivestringquantities, def:physics:phyx-mini-0166:phyxminiproblems-problemphyxmini0166-obeysstandingstringlaws}
  The pictured second harmonic has wavelength equal to the 0.500 m string length
\end{lemma}
\begin{proof}
  The conclusion follows from the governing relations and figure readouts named in the statement of wavelength In Meters eq one Half.
\end{proof}
\begin{lemma}[string Mass In Grams eq two Thousand Three Hundred Four div one Hundred Twenty Five]
  \label{lem:physics:phyx-mini-0166:phyxminiproblems-problemphyxmini0166-stringmassingrams-eq-twothousandthreehundredfour-div-onehundredtwentyfive}
  \lean{PhyXMiniProblems.ProblemPhyXMini0166.stringMassInGrams_eq_twoThousandThreeHundredFour_div_oneHundredTwentyFive}
  \uses{def:physics:phyx-mini-0166:phyxminiproblems-problemphyxmini0166-massingrams, def:physics:phyx-mini-0166:phyxminiproblems-problemphyxmini0166-vibratingstringsetup, def:physics:phyx-mini-0166:phyxminiproblems-problemphyxmini0166-matchesproblemreadouts, def:physics:phyx-mini-0166:phyxminiproblems-problemphyxmini0166-matchesprimaryfigure, def:physics:phyx-mini-0166:phyxminiproblems-problemphyxmini0166-haspositivestringquantities, def:physics:phyx-mini-0166:phyxminiproblems-problemphyxmini0166-obeysstandingstringlaws}
  The idealized exact inputs give 2304/125 g = 18.432 g. This exposes the unrounded physical calculation separately from selection among the displayed choices
\end{lemma}
\begin{proof}
  The conclusion follows from the governing relations and figure readouts named in the statement of string Mass In Grams eq two Thousand Three Hundred Four div one Hundred Twenty Five.
\end{proof}
% --- Archon declaration topology end ---
% --- Archon physics formalization source end ---
